\section{Methodology}
\label{sec:methodology}

The architecture of \textit{Cluster Around and Find Out} operates as an interactive pipeline that translates qualitative natural language conversations into quantitative geometric adjustments within a vector space. 
The workflow is divided into four main phases: initial corpus overview and baseline specification, prompt-engineered preference elicitation, instruction-guided embedding generation, and final customized clustering and evaluation pipeline. 
The system couples conversational efficiency tracking with multi-dimensional metric verification to evaluate end-to-end user alignment.

\subsection{Phase 1: Initial Corpus Overview and Baseline Specification}
The pipeline begins with an unlabelled text corpus consisting of a subset of 500 questions from the StackOverflow full dataset.
Before initializing the user interface, the system generates an initial structural layout of the text space to gather clean context for the conversational agent. 

To achieve this, the corpus is processed using a standard, conventional text embedding system. 
A traditional $K$-means clustering algorithm is executed over these conventional vectors with $K$ arbitrarily set to $10$. 
From each of the $10$ identified clusters, the system randomly samples $3$ representative documents. 
This yields a total of $30$ anchor documents that encapsulate the broad thematic variation of the dataset. 
These $30$ documents are fed into the system prompt of the conversational agent to ground its domain awareness.

Crucially, this phase also defines the uninstructed baseline used for downstream evaluation. 
For the baseline condition, the entire 500-document dataset is embedded using the primary encoder, \textbf{Qwen3-Embedding (8B)}, configured with a static instruction stating that the user has no specific sorting preferences.
The downstream clustering algorithm is then selected and configured by \textbf{Qwen3.6-35B-A3B} without any human alignment conditioning.
This baseline assignment is completely isolated from the conversation pipeline; 
it is never shown to the user or used by the interview agent, serving strictly as a control group to evaluate preference alignment post-experiment.

\subsection{Phase 2: Prompt-Engineered Preference Elicitation}
Instead of relying on heavy optimization or post-training adjustments, our conversational interface relies entirely on strategic prompt engineering. 
The system initializes the interview agent using \textbf{Qwen3.6-35B-A3B}. 
The model is configured via a comprehensive system prompt that embeds:
\begin{enumerate}
    \item The operational role of an expert data-mining interviewer.
    \item The dataset context derived from the Phase 1 sampling, passing the $30$ randomly sampled anchor documents directly into the prompt.
    \item Strict formatting rules corresponding to one of the three tested interactive paradigms: 
    \textit{yes/no}, \textit{multiple-choice}, or \textit{open-ended} question formats.
\end{enumerate}

During the dialogue loop, the agent interviews a simulated user configured with specific combinations of 3 personas and 6 goals.
The agent presents questions regarding document relationships, target granularity, and sorting logic, utilizing the initial cluster summaries to present concrete examples. 
The interaction continues dynamically until the conversation hits convergence criteria—either reaching a structural consensus or hitting an arbitrary max turn limit set to 20. 
At this point, the Qwen3.6-35B-A3B agent synthesizes the entire chat history into a concise, dense paragraph of text: 
the \textit{User Preference Profile}.

\subsection{Phase 3: Instruction-Guided Embedding Generation}
Once the User Preference Profile is synthesized, it is directly injected into the downstream embedding pipeline as a formatting instruction. 
We leverage \textbf{Qwen3-Embedding (8B)}, a decoder-only model optimized for instruction-tuned vector pooling. 

Unlike traditional static encoders that project identical text to identical vectors regardless of intent, the Qwen3-Embedding model accepts a structured instruction prefix. 
We construct our embedding prompt by mapping the User Preference Profile directly into the instruction field of the encoder. 
The raw StackOverflow questions are then passed to the model under this prefix. 
The model utilizes its internal bidirectional attention mechanisms to alter the weights of semantic tokens dynamically based on the instruction. 
This maps the text into an adjusted high-dimensional space where documents sharing user-preferred traits are projected closer together, and documents violating the user's criteria are pushed further apart.

\subsection{Phase 4: Final Customized Clustering and Evaluation Pipeline}
In the final phase, the system applies a clustering algorithm over the newly adjusted instruction-guided embedding space. 
Because the vector geometry itself has been altered by the Qwen3-Embedding model to reflect the conversation's nuances, a standard geometric partitioning algorithm can now capture subjective human intents. 

The resulting clusters represent the customized configuration tailored to the user's specific persona and goal. 
To maintain strict experimental alignment, the performance of these final clusters, along with the operational metrics of the dialogue loop, are captured, recorded, and processed using our LLM-as-judge. 
This allows for an end-to-end evaluation of how well the interactive prompting paradigm shifts the underlying mathematical structure of the data compared to our uninstructed Phase 1 control baseline.